\item Dos electrones en el mismo átomo tienen ambos $n = 3$ y $\ell = 1$. Suponga que los electrones son distinguibles.
\begin{enumerate}
    \item ¿Cuántos estados diferentes del átomo son posibles considerando los números cuánticos que estos dos electrones pueden tener?
    \item \textbf{¿Qué pasaría si?} ¿Cuántos estados serían posibles si el principio de exclusión de Pauli no fuera operativo?
\end{enumerate}